%% ==================== 模型假设 模板 (v1.5.6, BZD 数模社 2026 规范) ====================
%%
%% 借鉴 BZD 数模社 第三章 模板: 引导语 + 4-8 条假设 + 单条 "假设 + 说明"
%%
%% BZD 9 类假设:
%%  1. 数据真实性与有效性 (数据集完整且不存在错误, 实际是 "经清洗后保留的数据满足建模需要")
%%  2. 小概率事件排除 (战争/重大灾害等)
%%  3. 主要因素保留、次要因素忽略
%%  4. 系统平稳性 (环境/参数/状态在范围内稳定)
%%  5. 对象理想化 (质点/圆柱/均匀混合)
%%  6. 参数形式与概率分布 (服从某分布)
%%  7. 变量关系 (线性/可加/其他)
%%  8. 行为与决策 (决策者准则)
%%  9. 空间与时间 (同一周期口径一致)
%% 10. 独立性或无相互干扰
%%
%% 写法: 每条 "假设 + 说明 (依据/简化/作用)".
%% 关键假设建议 1-3 句话说明理由, 简单题目给的假设可简短.
%% 数量: 4-8 条最佳, 不硬性.
%%
%% 跑题用户必须:
%% 1. 删掉所有 % 注释 (整段 L1-L23)
%% 2. 替换【】占位符为自己的实际内容
%% 3. 删 下面的 \AssumptionSlot{...} 行 (check 15 抓这个 inline 【】)

\section{模型假设}

为突出研究对象的主要特征并降低次要因素对模型求解的干扰, 在不改变问题本质的前提下, 结合题目条件、数据特征及后续模型需要, 作出如下假设:

\begin{enumerate}[label=(\arabic*)]

\item \textbf{【假设 1 内容, 1 句。\textbf{粗体}突出关键术语】}。

\textbf{说明:} 【假设依据 + 简化对象 + 建模作用, 1-3 句。例如: "数据均来自题目附件, 经缺失值和异常值处理后保留的数据能够反映研究对象的基本特征。该假设是数据预处理和模型输入的合理性前提, 不影响后续建模结论。"】

\item \textbf{【假设 2 内容, 1 句】}。

\textbf{说明:} 【依据 + 简化对象 + 建模作用, 1-3 句】

\item \textbf{【假设 3 内容, 1 句】}。

\textbf{说明:} 【依据 + 简化对象 + 建模作用, 1-3 句】

\item \textbf{【假设 4 内容, 1 句】}。

\textbf{说明:} 【依据 + 简化对象 + 建模作用, 1-3 句】

\item \textbf{【假设 5 内容, 1 句】}。

\textbf{说明:} 【依据 + 简化对象 + 建模作用, 1-3 句】

\item \textbf{【假设 6 内容, 1 句 (可选)】}。

\textbf{说明:} 【依据 + 简化对象 + 建模作用, 1-3 句】

\end{enumerate}

% 不推荐开头: "为了方便模型建立, 使模型更加完备, 预测结果更加合理, 作出如下假设。" -- 空泛.
% 不可取: 假设模型必然正确 / 假设数据绝对没有缺失 / 假设后文不用的假设 / 重复近义句.

% 跑题后必须删下面这一行 (check 15 占位符自检)
\textbf{【这里粘贴从假设 1 到假设 6 (4-8 条) 拼起来的不超过 1 页内容, 删掉这一行】}
